\seccionanexo{Anexo D: Instrumentos de Diagnóstico del Proceso de Inscripción (Objetivo Específico 1)}

\newcommand{\casilla}{\framebox[0.38cm]{\rule{0pt}{0.33cm}}}

Instrumentos aplicados durante el primer bimestre de la gestión 2026 para diagnosticar las deficiencias operativas y los requerimientos funcionales del proceso de inscripción y matriculación del BTH, conforme al objetivo específico 1 y al diagnóstico 2022-2025 del Capítulo 2.

\noindent\textbf{Instrumento 1: Encuesta diagnóstica a estudiantes postulantes}

Cuestionario estructurado en escala tipo Likert de 5 puntos (1: Totalmente en desacuerdo a 5: Totalmente de acuerdo) aplicado a la muestra de 60 estudiantes postulantes del MÓDULO TECNOLÓGICO PRODUCTIVO SAN JULIÁN, de las siete especialidades y de las distintas unidades educativas de procedencia (anonimizado). Tiempo estimado: 6 min.

\noindent\textbf{Bloque A. Identificación (anonimizada):}

\begin{itemize}
    \item Unidad educativa de procedencia: \rule[-2pt]{3.5cm}{0.4pt}
    \item Especialidad de interés: \rule[-2pt]{3.5cm}{0.4pt}
    \item ¿Dispone de teléfono celular con acceso a internet? \casilla{} Siempre\quad \casilla{} A veces\quad \casilla{} Nunca
\end{itemize}

\noindent\textbf{Bloque B. Ítems Likert (1--5):}

\begin{enumerate}
    \item \textbf{Filas y espera presencial (2022-2024):}
    \begin{itemize}
        \item Durante el proceso de inscripción presencial tuve que hacer filas de varios días o permanecer de noche en el Módulo.
        \item La espera a la intemperie (sol, lluvia o frío) me obligó a dejar de asistir al colegio o a mis actividades de trabajo.
    \end{itemize}
    \item \textbf{Recuperación de credenciales (página de 2025):}
    \begin{itemize}
        \item En la gestión 2025 olvidé mi contraseña de la página de preinscripción.
        \item Para recuperar mi contraseña tuve que trasladarme presencialmente al Módulo y esperar varios días.
        \item Mientras esperaba el restablecimiento de mi contraseña, los cupos de la especialidad que deseaba se agotaron.
        \item Considero necesario recuperar la contraseña con mi cédula de identidad y fecha de nacimiento, sin asistir al Módulo y en menos de 2 minutos.
    \end{itemize}
    \item \textbf{Matriculación documental:}
    \begin{itemize}
        \item No sabía con exactitud qué documentos exigía el BTH para completar mi matriculación (libretas de 1ro a 3ro, cédula del tutor, carné y certificado de nacimiento).
        \item La entrega de mis documentos se realizó sin un registro claro y ordenado que me permitiera verificar su estado.
        \item Mi expediente quedó incompleto o se extravió algún documento durante el proceso de la gestión 2025.
        \item Considero fundamental que el sistema registre y verifique todos mis documentos y me entregue un comprobante con código de verificación.
    \end{itemize}
    \item \textbf{Distribución en cursos y paralelos:}
    \begin{itemize}
        \item Desconocía los criterios con los que se asignaría mi curso y paralelo (4to A o 4to B) según mi unidad educativa de procedencia.
        \item La asignación manual de paralelos generó cursos con cantidades desiguales de estudiantes o errores de registro.
        \item Presenté conflictos de horario entre el Módulo y mi unidad educativa de procedencia.
        \item Espero que mi asignación a curso y paralelo responda a criterios equitativos, transparentes y compatibles con mi horario escolar.
    \end{itemize}
    \item \textbf{Necesidades y requerimientos funcionales:}
    \begin{itemize}
        \item Prefiero realizar mi inscripción por internet desde mi teléfono celular, incluso con conexión de baja intensidad.
        \item La espera y los gastos de traslado (pasajes y fotocopias) afectaron económicamente a mi familia durante el proceso de inscripción.
    \end{itemize}
\end{enumerate}

\noindent\textbf{Bloque C. Preguntas cerradas:}

\begin{itemize}
    \item Tiempo total de espera presencial estimado en la gestión 2025: \casilla{} Menos de 1 día\quad \casilla{} 1 a 2 días\quad \casilla{} 3 días o más
    \item Gasto estimado en pasajes y fotocopias durante el proceso (Bs): \casilla{} Menos de 20\quad \casilla{} 20 a 50\quad \casilla{} Más de 50
    \item Mayor dificultad encontrada: \casilla{} Recuperar la contraseña\quad \casilla{} Reunir y entregar los documentos\quad \casilla{} Obtener la especialidad y paralelo deseados
    \item Forma preferida de recuperar la contraseña: \casilla{} Por celular con cédula y fecha de nacimiento\quad \casilla{} Asistiendo al Módulo\quad \casilla{} Correo electrónico\quad \casilla{} Otra: \rule[-2pt]{2.5cm}{0.4pt}
    \item Requerimiento más importante para la página de inscripción: \casilla{} Recuperación autónoma de la contraseña\quad \casilla{} Registro y verificación de documentos\quad \casilla{} Asignación automática y transparente de paralelos\quad \casilla{} Comprobante descargable con código de verificación
\end{itemize}

\vspace{0.4cm}
\noindent\textbf{Instrumento 2: Guía de Entrevista Semiestructurada — Personal Administrativo y Docente (5 funcionarios: Director, Secretaría Académica y 3 docentes entre ellos Tutores de Curso):}

\begin{enumerate}
    \item ¿Qué dificultades enfrentaron en la gestión 2025 con los reseteos manuales de contraseñas y la pérdida de cupos vocacionales?
    \item ¿Cómo se recepcionaban y archivaban las libretas de 1ro a 3ro de secundaria, cédulas de tutores, carné y certificado antes de la V2.0?
    \item ¿Qué porcentaje de expedientes quedaba incompleto o sin trazabilidad en 2025?
    \item ¿Qué errores y sobrecarga generaba la asignación manual a paralelos en planillas externas por unidad educativa?
    \item ¿Qué conflictos de coordinación horaria surgían entre el MTP y las unidades educativas de origen?
    \item ¿Qué requisitos de validación consideran imprescindibles para el módulo de matriculación documental auditada?
    \item ¿Qué criterios de equilibrio por taller deberían aplicarse en la asignación automatizada (4to A/B)?
    \item ¿Qué expectativas tienen sobre el panel de control y reportes en tiempo real para la toma de decisiones?
    \item ¿Cuántos restablecimientos manuales de contraseña se atendieron por gestión en ventanilla?
    \item ¿Cuántas horas semanales dedica su área al proceso de inscripción y matriculación?
    \item ¿Qué porcentaje de postulantes terminó la gestión 2025 en una especialidad distinta a la que deseaba?
\end{enumerate}

\notatabla{Instrumentos de diagnóstico del objetivo específico 1, aplicados en el primer bimestre de 2026 a 60 estudiantes postulantes y 5 funcionarios del MTP San Julián (muestreo por conveniencia). Datos anonimizados (C.I. como ****). Validación por juicio del tutor y prueba piloto con 5 estudiantes (ver Anexo B). Los ítems se alinean con las tres brechas del Capítulo 2: recuperación de credenciales, matriculación documental y distribución por unidad educativa.}
